%============================================================
% BAB 2: GAMBARAN PERUSAHAAN
%   2.1 Profil Perusahaan (termasuk visi & misi dalam narasi)
%   2.2 Struktur Organisasi Perusahaan
%       2.2.x per divisi
%   2.3 Hak dan Wewenang
%       2.3.x per jabatan
%   2.4 Lokasi Perusahaan
%   2.5 Etika Profesi
%============================================================
\chapter{GAMBARAN PERUSAHAAN}

\section{PROFIL PERUSAHAAN}

Jelaskan profil singkat \NamaPerusahaan, meliputi: tahun berdiri, bidang usaha
utama, skala perusahaan, dan gambaran umum posisinya di industri atau sektor
yang bersangkutan.

Sertakan juga visi dan misi perusahaan dalam bagian ini sebagai narasi
atau daftar ringkas. Jelaskan nilai-nilai utama yang dijalankan perusahaan.

% Contoh penyisipan logo:
% \begin{figure}[H]
%   \centering
%   \includegraphics[width=4cm]{assets/gambar/logo_perusahaan.png}
%   \caption{Logo \NamaPerusahaan}
%   \label{fig:logo_perusahaan}
% \end{figure}

\section{STRUKTUR ORGANISASI PERUSAHAAN}

Struktur organisasi \NamaPerusahaan\ dirancang untuk mendukung efisiensi
operasional dan kolaborasi antar departemen. Berikut adalah struktur organisasi
\NamaPerusahaan.

% Contoh penyisipan bagan organisasi:
% \begin{figure}[H]
%   \centering
%   \includegraphics[width=0.9\textwidth]{assets/gambar/struktur_organisasi.png}
%   \caption{Struktur Organisasi \NamaPerusahaan}
%   \label{fig:struktur_org}
% \end{figure}

% Sesuaikan sub-section dengan divisi yang ada di perusahaan.
\subsection{[Nama Divisi 1]}

\begin{itemize}[leftmargin=*, itemsep=0pt]
  \item \textbf{Kepala Divisi:} Nama kepala divisi.
  \item \textbf{Tugas:} Deskripsi tugas dan tanggung jawab divisi ini.
\end{itemize}

\subsection{[Nama Divisi 2]}

\begin{itemize}[leftmargin=*, itemsep=0pt]
  \item \textbf{Kepala Divisi:} Nama kepala divisi.
  \item \textbf{Tugas:} Deskripsi tugas dan tanggung jawab divisi ini.
\end{itemize}

\section{HAK DAN WEWENANG}

Setiap unit kerja di \NamaPerusahaan\ memiliki hak, wewenang, serta tugas yang
jelas untuk memastikan operasional perusahaan berjalan lancar dan efisien.

% Sesuaikan sub-section dengan jabatan struktural yang ada.
\subsection{[Jabatan Utama 1]}

\begin{itemize}[leftmargin=*, itemsep=0pt]
  \item \textbf{Hak:} Hak jabatan ini.
  \item \textbf{Wewenang:} Wewenang jabatan ini.
  \item \textbf{Tugas:} Tugas utama jabatan ini.
\end{itemize}

\subsection{[Jabatan Utama 2]}

\begin{itemize}[leftmargin=*, itemsep=0pt]
  \item \textbf{Hak:} Hak jabatan ini.
  \item \textbf{Wewenang:} Wewenang jabatan ini.
  \item \textbf{Tugas:} Tugas utama jabatan ini.
\end{itemize}

\section{LOKASI \MakeUppercase{\NamaPerusahaan}}

\NamaPerusahaan\ berlokasi di \AlamatPerusahaan. Lokasi ini strategis dan mudah
diakses, mendukung operasional perusahaan serta memudahkan kolaborasi dengan
klien dan mitra bisnis.

\begin{itemize}[leftmargin=*, itemsep=0pt]
  \item \textbf{Nama Perusahaan:} \NamaPerusahaan
  \item \textbf{Alamat:} \AlamatPerusahaan
\end{itemize}

% Contoh penyisipan peta:
% \begin{figure}[H]
%   \centering
%   \includegraphics[width=0.85\textwidth]{assets/gambar/peta_lokasi.png}
%   \caption{Peta Lokasi \NamaPerusahaan}
%   \label{fig:peta_lokasi}
% \end{figure}

\section{ETIKA PROFESI}

Uraikan etika dan budaya kerja yang diterapkan di \NamaPerusahaan, meliputi
standar perilaku, kedisiplinan, kerja sama tim, komunikasi, dan nilai-nilai
profesionalisme yang berlaku di lingkungan kerja instansi tersebut.

\begin{enumerate}[label=\arabic*., leftmargin=*, itemsep=0pt]
  \item \textbf{Sikap dan Kedisiplinan:} Deskripsi etika pertama.
  \item \textbf{Kerjasama Tim:} Deskripsi etika kedua.
  \item \textbf{Komunikasi yang Efektif:} Deskripsi etika ketiga.
  \item \textbf{Integritas:} Deskripsi etika keempat.
\end{enumerate}
